\scanpage{10}
\mspage{-0.15-}
served under flat base extension -- hence the analogous sequence $(*)^{\handwritten{\prime}}$ on $Y^{\handwritten{\prime}}$ is exact.

(8). On the other hand, consider the converse. I claim that if $Y^{\handwritten{\prime}}\handwritten{\longrightarrow} Y$ is faithfully flat and quasi-compact, and if the extended morphism $\handwritten{\phi}^{\handwritten{\prime}}:X^{\handwritten{\prime}}\handwritten{\longrightarrow} Y^{\handwritten{\prime}}$ makes $(Y^{\handwritten{\prime}},\handwritten{\phi}^{\handwritten{\prime}})$ a universal categorical quotient or geometric quotient or universal geometric quotient of $X^{\handwritten{\prime}}$ by $G$, then $(Y,\handwritten{\phi})$ is the same type of quotient of $X$ by $G$. These assertions are all simple corollaries of the general theory of descent of SGA 8.

\fbox{\S3.} We shall now look somewhat more closely at the structure of an action of a group. As above, let $G/S$ act via $\handwritten{\sigma}$ on $X/S$.

\uline{Definition} \uline{0.8}: The action $\handwritten{\sigma}$ is said to be

i) \uline{closed} if for all geometric points $x$ of $X$, the orbit $O(x)\handwritten{\subset}\overline X$ is closed, (i.e. $\overline X=X\times_S\operatorname{Spec}\handwritten{\Omega}$ if $x$ is an $\handwritten{\Omega}$-valued point),

ii) \uline{separated} if the image of
\[
\handwritten{\unsure{\Psi|\overline{\underline\psi}}}=(\handwritten{\sigma},p_2):G\times_S X\handwritten{\longrightarrow} X\times_S X
\]
is closed,

iii) \uline{proper} if $\handwritten{\Psi}$ is proper,

iv) \uline{free} if $\handwritten{\Psi}$ is \tikz[baseline=(v.base)]{\node[draw,ellipse,inner xsep=2pt,inner ysep=1pt] (v) {a closed immersion.};}
% slanted
\margin{\handwritten{\fbox{monomorphisme}}}
% slanted
\margin{\handwritten{conflits avec des terminologies reçues dire : \uline{proper and free}}}

Moreover, let $\handwritten{\mathscr S}=S(1_X)$ be the stabilizer of $1_X:X\handwritten{\longrightarrow} X$ (cf. Definition 0.4). It is a sub-group pre-scheme of the group pre-scheme $G\times_S X/X$. Let $\handwritten{\overline\omega}:\handwritten{\mathscr S}\handwritten{\longrightarrow} X$ be the projection and let $e_X:X\handwritten{\longrightarrow}\handwritten{\mathscr S}$

% archive sheet number: 9
